% Generated draft: requires mathematical and source review.
\documentclass[11pt]{article}
\usepackage[margin=25mm]{geometry}
\usepackage{fontspec}
\setmainfont{Latin Modern Roman}
\usepackage{amsmath,amssymb,mathtools,unicode-math}
\setmathfont{Latin Modern Math}
\usepackage{enumitem,needspace,xurl,hyperref}
\setlength{\parindent}{0pt}
\setlength{\parskip}{5pt}
\setlength{\emergencystretch}{3em}
\newcommand{\N}{\mathbb{N}}
\newcommand{\Z}{\mathbb{Z}}
\newcommand{\Q}{\mathbb{Q}}
\newcommand{\R}{\mathbb{R}}
\newcommand{\C}{\mathbb{C}}
\newcommand{\F}{\mathbb{F}}
\newcommand{\A}{\mathbb{A}}
\newcommand{\PP}{\mathbb{P}}
\newcommand{\E}{\mathbb{E}}
\newcommand{\eps}{\varepsilon}
\newcommand{\dd}{\,\mathrm{d}}
\begin{document}
\section*{UnsolvedMath ID: 30002208}
\textbf{Source problem code:} catalogue entry
\subsection{Definitions and mathematical statement}
For this draft, use the following explicit conventions. The supplied record does not define its largeness terminology, so agreement with the source, particularly for piecewise Bohr sets, requires review.

For subsets \(A,B\subseteq\Z\), define
\[
A+B=\{a+b:a\in A,\ b\in B\}.
\]
The upper Banach density of \(B\subseteq\Z\) is
\[
d^*(B)=\limsup_{N\to\infty}\ \sup_{m\in\Z}
\frac{\lvert B\cap\{m+1,\ldots,m+N\}\rvert}{N}.
\]
A set \(T\subseteq\Z\) is thick if, for every integer \(N\geq1\), there exists \(m\in\Z\) such that \(\{m+1,\ldots,m+N\}\subseteq T\). A set \(S\subseteq\Z\) is syndetic if \(F+S=\Z\) for some finite nonempty \(F\subseteq\Z\). It is piecewise syndetic if \(F+S\) is thick for some finite nonempty \(F\subseteq\Z\).

Write \(\lVert x\rVert_{\R/\Z}=\inf_{k\in\Z}|x-k|\). A basic Bohr neighborhood of zero is a set
\[
V(\alpha_1,\ldots,\alpha_r;\varepsilon)
=\{n\in\Z:\lVert n\alpha_j\rVert_{\R/\Z}<\varepsilon\text{ for every }1\leq j\leq r\},
\]
where \(r\geq1\), \(\alpha_1,\ldots,\alpha_r\in\R\), and \(\varepsilon>0\). Here a Bohr set means a subset \(U\subseteq\Z\) containing a translate \(k+V(\alpha_1,\ldots,\alpha_r;\varepsilon)\) for some such parameters and \(k\in\Z\). A set \(S\subseteq\Z\) is piecewise Bohr if there exist a Bohr set \(U\) and a thick set \(T\) such that \(U\cap T\subseteq S\).

The source asks what can be said about a set \(A\subseteq\Z\) satisfying either of the following alternative hypotheses, considered separately:
\[
\begin{aligned}
\mathcal P_{\mathrm{ps}}(A):\quad&
\text{for every }B\subseteq\Z,\quad
 d^*(B)>0\ \Longrightarrow\ A+B\text{ is piecewise syndetic};\\
\mathcal P_{\mathrm{pB}}(A):\quad&
\text{for every }B\subseteq\Z,\quad
 d^*(B)>0\ \Longrightarrow\ A+B\text{ is piecewise Bohr}.
\end{aligned}
\]
Thus the task is to describe or characterize the sets \(A\) satisfying each hypothesis. The source does not specify a proposed conclusion, a particular characterization, or a precise conjectural equivalence. These predicates formalize the supplied hypotheses without supplying a missing conclusion.

The supplied dated literature assessment says that the extracted question does not define \(A\). This conflicts with the supplied original and cleaned statements, which explicitly give \(A\subseteq\Z\) and the hypotheses above. The intended scope of the requested description remains unresolved.
\subsection{Short English statement}
Describe the subsets \(A\subseteq\Z\) whose sum with every set of positive upper Banach density is piecewise syndetic. Consider separately the alternative requirement that every such sum be piecewise Bohr. The source asks a broad characterization question and supplies no specific proposed conclusion; its terminology and conflicting curation notes require review.
\subsection{Sources}
\begin{itemize}
\item Supplied local UnsolvedMath registry record, permanent ID \(30002208\), source problem code \texttt{catalogue entry}, titled \emph{Additive Sets Preserving Piecewise Syndeticity}, in the Oberwolfach Reports Open Problems collection. The record supplies original and cleaned statements; no individual proposer is identified.
\item Supplied report citation: \emph{Arbeitsgemeinschaft: Ergodic Theory and Combinatorial Number Theory} (2013), Oberwolfach Reports, pp.~2985--3059, DOI \texttt{10.4171/owr/2012/50}. Supplied source URL and literature-source link labeled ``original OWR report'': \url{https://doi.org/10.4171/owr/2012/50}. The background describes these reports as documenting mathematical workshops and their open-problem discussions. The record separately gives proposed year 2012.
\item Literature assessment dated exactly \texttt{2026-08-21}: the supplied curation labels the problem open and says that the extraction is only the question ``What can be said about A?'', does not define \(A\), and cannot be classified more narrowly. This is supplied dated curation, not an independently verified literature assessment. It conflicts with the fuller statements supplied in the same record. The background requests independent verification of time-sensitive status and cited claims.
\item The supplied statement-status label is \texttt{corrected\_verified}; its accompanying note reports manual semantic, typographic, and LaTeX review and calls the cleanup self-contained and meaning-preserving. These are supplied registry claims, not verification performed for this draft. The cited report has not been consulted for this draft.
\end{itemize}
\end{document}
